\chapter{Hemoglobin A1c (HbA1c)}

\section*{What Is This Test?}
The hemoglobin A1c (HbA1c) test measures the percentage of hemoglobin that has been irreversibly glycated --- the fraction of the A1c subform (hemoglobin A1c) in which glucose is non-enzymatically attached to the N-terminal valine of the beta-globin chain. Because red blood cells live approximately 120 days and glucose attaches slowly and continuously, the HbA1c percentage reflects the average blood glucose level over the preceding 8--12 weeks (weighted toward the most recent 4--8 weeks). HbA1c is the cornerstone test for both the diagnosis of diabetes mellitus and the long-term monitoring of glycemic control. It is the only routine diabetes test that provides an integrated, retrospective measure of average glycemia without requiring the patient to fast or to log glucose values. Each 1\% increase in HbA1c corresponds to an average glucose increase of approximately 28--29 mg/dL (1.6 mmol/mol per mmol/mol in IFCC units). Because chronic hyperglycemia drives the microvascular and macrovascular complications of diabetes, HbA1c is the single best predictor of diabetes-related complications in population studies and in individual patient management.

\section*{Alternative Names}
A1c; Glycated Hemoglobin; Glycohemoglobin; Glycosylated Hemoglobin; Hemoglobin A1c; HbA1C; Glycated Hemoglobin (GHb); Total Glycated Hemoglobin (distinct from glycated albumin or fructosamine)

\section*{Principle}
HbA1c forms through a slow, non-enzymatic condensation between glucose and the N-terminal amino group of the beta-globin valine, a reaction that proceeds continuously throughout the red cell's 120-day lifespan. The rate of formation is directly proportional to the ambient glucose concentration, so the steady-state HbA1c reflects the integrated average glucose over the preceding 2--3 months. Laboratory methods fall into four groups: (1) High-performance liquid chromatography (HPLC) using cation-exchange columns, which physically separates glycated and non-glycated hemoglobin fractions and is the gold standard (ngsp-certified); (2) Immunoassay methods using antibodies specific for the glycated N-terminal sequence of the beta chain; (3) Capillary electrophoresis, which separates hemoglobin species by charge; (4) Enzymatic methods that quantify glycated beta-globin N-terminal residues. Results are standardized to the National Glycohemoglobin Standardization Program (NGSP) against the Diabetes Control and Complications Trial (DCCT) reference method and reported as a percentage; the International Federation of Clinical Chemistry (IFCC) reference method reports mmol/mol, with the conversion: NGSP \% = 0.09148 $\times$ IFCC (mmol/mol) + 2.152. Because HbA1c is a fraction of total hemoglobin, results are only valid when hemoglobin is normal and red cell survival is normal; conditions that alter red cell lifespan or introduce abnormal hemoglobins distort the value.

\section*{History}
The glycated hemoglobin was discovered by Samuel Rahbar in 1968, who noticed an unusual fast-moving hemoglobin fraction on electrophoresis in patients with diabetes. In 1971 Trivelli and colleagues confirmed that this fraction was elevated in diabetes and represented non-enzymatic glycosylation of hemoglobin. The 1976 work of Koenig and Cerami established that the amount of glycosylated hemoglobin reflected integrated blood glucose over the prior weeks, giving rise to the concept of a ``memory'' of glycemia. The landmark Diabetes Control and Complications Trial (DCCT, published 1993) and the United Kingdom Prospective Diabetes Study (UKPDS, published 1998) demonstrated that lowering HbA1c reduces microvascular complications (retinopathy, nephropathy, neuropathy) and led to the universal adoption of HbA1c targets and standardized reporting. In 2002 the NGSP harmonized assay standardization worldwide. In 2010 the American Diabetes Association (ADA) formally added HbA1c $\geq$6.5\% to fasting glucose and OGTT as a diagnostic criterion for diabetes, a position endorsed by the World Health Organization in 2011.

\section*{Why This Test Is Ordered}
\begin{itemize}
  \item \textbf{Screening for prediabetes and diabetes in asymptomatic adults:} All adults aged 35 years and older with risk factors, and all adults aged 45 years and older, regardless of risk factors; more frequent testing in those with prediabetes or strong risk factors
  \item \textbf{Diagnosis of diabetes mellitus:} HbA1c $\geq$6.5\% on two separate occasions confirms the diagnosis (a single value $\geq$6.5\% with symptoms of hyperglycemia may suffice)
  \item \textbf{Monitoring glycemic control in established diabetes:} Recommended every 3 months when treatment is being adjusted or glucose is uncontrolled, and every 6 months in stable patients meeting targets
  \item \textbf{Assessment of overall glycemic control in patients on insulin or oral agents} to guide therapy intensification
  \item \textbf{Evaluation during pregnancy in women with pre-existing diabetes} (not for gestational diabetes screening, where it is less sensitive)
  \item \textbf{Prediction of complication risk:} Higher HbA1c correlates with higher risk of retinopathy, nephropathy, and neuropathy
\end{itemize}

\section*{Is This a Primary or Secondary Test}
HbA1c is a primary test for the diagnosis and longitudinal monitoring of diabetes mellitus. When an asymptomatic result is $\geq$6.5\%, it must be repeated on a second sample (with rare exceptions) for confirmation. When used for monitoring, it is the primary outcome measure of glycemic control but is always interpreted together with self-monitored blood glucose values and, in some settings, with glycated albumin (fructosamine).

\section*{How This Test Is Performed}
A venous blood sample is collected in a lavender-top (EDTA) tube, or a small capillary sample may be taken by fingerstick for point-of-care devices (which require NGSP certification and correlation with laboratory methods). Approximately 2--3 mL of whole blood is sufficient. The sample is analyzed on an automated HPLC, immunoassay, capillary electrophoresis, or enzymatic platform. Fasting is not required because HbA1c is unaffected by short-term glucose excursions. Point-of-care A1c testing provides results within minutes during the office visit; laboratory testing typically returns results within 1--2 days. The venipuncture itself takes less than 5 minutes.

\section*{What Preparation Is Needed}
No fasting or special preparation is required. The test can be performed at any time of day. The patient should inform the healthcare provider of all medications (especially erythropoiesis-stimulating agents, iron supplements, and any recent blood transfusion), known hemoglobinopathies, chronic kidney disease, pregnancy, and any history of anemia or recent blood loss, because these conditions alter red blood cell turnover and can make HbA1c unreliable.

\section*{Contraindications and Precautions}
There are no absolute contraindications to the blood draw. However, HbA1c is \textbf{not reliable} when red blood cell survival is abnormal or when abnormal hemoglobins are present. Conditions that \textbf{falsely lower} HbA1c include hemolytic anemias (shortened red cell survival), recent significant blood loss or transfusion, pregnancy (dilutional and shortened red cell survival), and chronic kidney disease treated with erythropoiesis-stimulating agents. Conditions that \textbf{falsely raise} HbA1c include iron deficiency anemia, vitamin B12 or folate deficiency, asplenia, and uremia (carbamylated hemoglobin may interfere with some methods). Hemoglobin variants (HbS, HbC, HbE, HbD, and HbF) interfere with many immunoassay and HPLC methods; the clinical laboratory must be informed of a known variant so that an appropriate variant-robust method (or alternative monitoring with glycated albumin or glucose management indicator) is used. In patients with sickle cell trait, renal failure, pregnancy, or known variants, fructosamine or glycated albumin should be used instead. In people of African, Mediterranean, or Southeast Asian descent, discordance between glucose and HbA1c should prompt evaluation for a hemoglobin variant or red cell disorder.

\section*{Risks}
Risks are those of routine venipuncture: mild pain, bruising, small hematoma at the collection site, lightheadedness or vasovagal syncope, and extremely rarely, infection or nerve injury. No clinically significant risk is associated with the assay itself.

\section*{What Happens After the Test}
Point-of-care results are discussed at the same visit; laboratory results are usually available within 1--2 days. If the result confirms diabetes ($\geq$6.5\%), the healthcare provider will review lifestyle measures (diet, exercise, weight loss) and, in most cases, prescribe metformin along with cardiovascular risk management (blood pressure, lipids, smoking cessation, weight). If the result indicates prediabetes (5.7--6.4\%), intensive lifestyle intervention is recommended to prevent progression, and the test is repeated annually (or every 1--2 years after normalization). For patients with established diabetes, the HbA1c result is compared against the individualized target (typically $<$7\% for most nonpregnant adults, $<$6.5\% for younger patients with new disease and no significant cardiovascular disease if it can be achieved without significant hypoglycemia, and $<$8\% for frail elderly patients), and therapy is adjusted accordingly. For patients on insulin, the A1c trend is used together with self-monitored glucose and, when available, continuous glucose monitoring data (which can provide the glucose management indicator, GMI).

\section*{Understanding Results}

\subsection*{Below Normal (Less than 5.7\%)}
\begin{itemize}
  \item \textbf{Normal glucose tolerance:} HbA1c $<$5.7\% with normal fasting and postprandial glucose indicates normal glycemic status
  \item \textbf{Hemolytic anemia or recent blood loss:} Falsely low HbA1c due to shortened red cell survival
  \item \textbf{Recent blood transfusion:} Donor red cells dilute the patient's glycated fraction
  \item \textbf{Pregnancy:} Physiological hemodilution and reduced red cell lifespan lower HbA1c in the second and third trimesters
  \item \textbf{Erythropoiesis-stimulating agents or recent iron therapy:} Young red cells have had less time to accumulate glycation
\end{itemize}

\subsection*{Above Normal (5.7\% and Higher)}
\begin{itemize}
  \item \textbf{Prediabetes (5.7--6.4\%):} Increased risk of progression to diabetes; intensive lifestyle intervention reduces progression by approximately 58\%
  \item \textbf{Diabetes mellitus ($\geq$6.5\%):} Established diabetes; requires confirmation on repeat testing when asymptomatic
  \item \textbf{Poorly controlled diabetes on monitoring:} HbA1c above the individualized target indicates a need to intensify therapy; a change of 1\% corresponds to an average glucose change of approximately 28--29 mg/dL
  \item \textbf{Iron deficiency anemia, B12 or folate deficiency:} Falsely elevated HbA1c with improved red cell survival or increased immature cells
  \item \textbf{Chronic kidney disease with uremia:} Carbamylated hemoglobin may interfere with some methods, falsely elevating the result
\end{itemize}

\section*{Normal Results}
\begin{itemize}
  \item \textbf{Normal (no diabetes):} $<$5.7\% (39 mmol/mol or less in IFCC units)
  \item \textbf{Prediabetes:} 5.7--6.4\% (39--46 mmol/mol)
  \item \textbf{Diabetes:} $\geq$6.5\% (48 mmol/mol or higher)
  \item \textbf{Treatment target (most nonpregnant adults):} $<$7.0\% (53 mmol/mol)
  \item \textbf{More stringent target (selected younger patients):} $<$6.5\% if achievable without significant hypoglycemia
  \item \textbf{Less stringent target (frail or complex patients):} $<$8.0\% (64 mmol/mol)
\end{itemize}

\section*{Follow-up Tests}
\begin{itemize}
  \item \textbf{Repeat HbA1c:} For confirmation of a diagnostic result, and every 3--6 months for monitoring
  \item \textbf{Fasting plasma glucose:} In conjunction with diagnosis and when discordance with HbA1c requires resolution
  \item \textbf{Oral glucose tolerance test (OGTT):} When HbA1c is borderline or discordant, or for gestational diabetes screening
  \item \textbf{Urine albumin-to-creatinine ratio and serum creatinine:} Screening for diabetic kidney disease
  \item \textbf{Lipid panel:} Cardiovascular risk assessment, performed at diagnosis and periodically
  \item \textbf{Liver function tests:} Baseline before statin therapy and for monitoring of certain glucose-lowering medications
  \item \textbf{Fructosamine or glycated albumin:} When HbA1c is unreliable (hemoglobinopathy, hemolysis, pregnancy, recent transfusion, advanced kidney disease)
  \item \textbf{Continuous glucose monitoring:} To provide the glucose management indicator (GMI) and detailed glucose patterns in insulin-treated patients
\end{itemize}

\section*{References}
\begin{enumerate}
  \item American Diabetes Association. Standards of Care in Diabetes --- 2023. Diabetes Care. 2023;46(Suppl 1):S1--S291.
  \item The Diabetes Control and Complications Trial Research Group. The effect of intensive treatment of diabetes on the development and progression of long-term complications in insulin-dependent diabetes mellitus. N Engl J Med. 1993;329(14):977--986.
  \item UK Prospective Diabetes Study (UKPDS) Group. Intensive blood-glucose control with sulphonylureas or insulin compared with conventional treatment and risk of complications in patients with type 2 diabetes (UKPDS 33). Lancet. 1998;352(9131):837--853.
  \item Rahbar S. An abnormal hemoglobin in red cells of diabetics. Clin Chim Acta. 1968;22(2):296--298.
  \item World Health Organization. Use of glycated haemoglobin (HbA1c) in the diagnosis of diabetes mellitus. Diabetes Res Clin Pract. 2011;93(3):299--309.
  \item Burtis CA, Ashwood ER, Bruns DE. Tietz Textbook of Clinical Chemistry and Molecular Diagnostics. 7th ed. Elsevier; 2023.
  \item Fischbach FT, Dunning MB. A Manual of Laboratory and Diagnostic Tests. 11th ed. Wolters Kluwer; 2023.
  \item Little RR, Rohlfing CL, Sacks DB. Status of hemoglobin A1c measurement and goals for improvement: from chaos to order for improving diabetes care. Clin Chem. 2011;57(2):205--214.
\end{enumerate}

\clearpage
